\chapter{Methodology}

\section{Research Approach}

This study adopts a qualitative research approach based on a literature review. Rather than collecting primary data, the research focuses on analysing existing academic studies related to artificial intelligence (AI), generative AI, and consumer perception in marketing.

\section{Research Design}

The study uses thematic analysis to identify patterns and themes across selected academic literature. This approach allows the research to organise and interpret complex findings from multiple sources.

\section{Data Collection}

The data for this study consists of academic journal articles related to AI in marketing and consumer perception. These articles were selected using academic databases such as Google Scholar, Scopus, and ScienceDirect.

Keywords used for searching include:

\begin{itemize}
    \item Artificial intelligence in marketing
    \item Generative AI advertising
    \item Consumer perception AI
    \item AI-generated content
    \item AI trust, authenticity, and privacy
\end{itemize}

\section{Data Selection Criteria}

To ensure relevance and quality, the following criteria were applied:

\begin{itemize}
    \item Peer-reviewed journal articles
    \item Published between 2020 and 2025
    \item Focus on AI or generative AI in marketing
    \item Discuss consumer perception or related concepts
\end{itemize}

A total of approximately 15–25 articles were selected for detailed analysis.

\section{Thematic Analysis}

Thematic analysis was used to identify key themes across the selected literature. The process followed several steps:

\begin{enumerate}
    \item Familiarisation with the data through reading the articles
    \item Initial coding of key concepts and findings
    \item Grouping codes into broader themes
    \item Reviewing and refining themes
\end{enumerate}

Key themes identified include trust, authenticity, personalisation, privacy concerns, and purchase intention.

\section{Limitations}

This study has several limitations. First, it relies only on secondary data, which may limit the depth of analysis. Second, the selection of literature may introduce bias. Third, the findings are based on interpretation and may not represent all consumer perspectives.

\section{Summary}

This chapter has outlined the research approach, data collection process, and analytical method used in this study. The next chapter presents the findings based on thematic analysis.